\chapter{Inequalities}

\emph{Inequalities} play a big role in all of mathematics, but especially at the start of real analysis it is 
of major importance to get familiar with them. 

\section{The Absolute Value}

The \emph{absolute value} of a real number \(x\), denoted by \(|x|\), is defined as
\[
    |x| = 
    \begin{cases} 
        x & \text{if } x \ge 0, \\
        -x & \text{if } x < 0.
    \end{cases}
\]

\subsection{Properties of Absolute Value}

\begin{itemize}

    \item \(|x| \ge 0\)

    \item \(|x| = 0 \iff x = 0\)

    \item \(|-x| = |x|\)

    \item \(|xy| = |x| \cdot |y|\)

    \item \(\left|\frac{x}{y}\right| = \frac{|x|}{|y|}, \quad y \ne 0\)

\end{itemize}

\subsection{Interval Property of the Absolute Value}

Let \(a, b, x \in \R\) and assume that \(b \ge 0\). 

\[
    |x - a| \le b \implies a - b \le x \le a + b
\]

\textbf{Proof:}

Assume that \(|x - a| \le b\). Now we have two cases. 

If \(x - a \ge 0\), then \(0 \le x - a \le b\), which implies \(a \le x \le a + b\). Hence,

\[
   a - b \le x \le a + b.
\]

If \(x - a < 0\), then \(|x - a| = -(x - a) = -x + a\), which implies \(0 \le -x + a \le b\), and hence \(a - b \le x \le a\). Therefore, combining both cases, we get the desired interval property.

\[
   a - b \le x \le a + b.
\]

\QED

\section{Triangle Inequality}

Let \(x, y \in \R\). The following inequalities hold:

\begin{itemize}

    \item \(|x + y| \le |x| + |y|\)

    \item \(|x - y| \ge |x| - |y|\)

    \item \(||x| - |y|| \le |x - y|\)

    \item \(xy \le |x| \cdot |y|\)

\end{itemize}

\textbf{Proof of the triangle inequality:}

Let \(x \ge 0\) and \(y \ge 0\), then \(|x + y| = x + y\) and \(|x| = x, |y| = y\). Therefore,

\[
    |x + y| = x + y = |x| + |y|.
\]

If \(x < 0\) and \(y < 0\), then \(|x + y| = -(x + y) = -x - y\) and \(|x| = -x, |y| = -y\). Therefore,

\[
    |x + y| = -x - y = |x| + |y|.
\]

If \(x \ge 0\) and \(y < 0\), then \(|x + y| = x + (-y) = x - y\) and \(|x| = x, -|y| = y\). Because 

\[
    -|y| < |y| \implies x - y < x + |y| = |x| + |y|.
\]

Similarly, for the case \(x < 0\) and \(y \ge 0\), then \(|x + y| = -x + y\) and \(|x| = -x, |y| = y\). Because 

\[
    -|x| < |x| \implies -x + y < |x| + |y|.
\]

\QED

\section{Reciprocal Inequality}

\[
    a < b < c \implies \frac{1}{a} > \frac{1}{b} > \frac{1}{c}
\]

\section{Bernoulli Inequality}

Given \(x \in \R, x \ge -1, x \ne 0\), then for \(n \in \N, n \ge 1\), we have

\[
    (1 + x)^n \ge 1 + nx
\]

This inequality is known as the \emph{Bernoulli inequality} and it is useful for proving that some functions are stricly greater than other 
on some intervals of the real line, among other applications.

\textbf{Proof:}

Suppose \(a \ge -1\). 

Let our \emph{base case} be \(n = 1\). Then \((1 + x)^1 = 1 + x\), which satisfies the inequality \(1 + x \ge 1 + x\).

\emph{Induction hypothesis}: suppose the inequality holds for some \(n = k\), i.e.,

\[
    (1 + x)^k \ge 1 + kx
\]

\emph{Induction Step}: We want to prove that \((1 + x)^{k + 1} \ge 1 + (k  + 1)x\).

\begin{align*}
    (1 + x)^{k + 1} &= (1 + x)^k (1 + x) \\ 
    &\ge (1 + xn)(1 + x) \\ 
    &= 1 + x + xn + nx^2 \\ 
    &= 1 + (k + 1)x + nx^2 \\ 
    &\ge 1 + (k + 1)x
\end{align*}

In this induction proof we could not just use the induction hypothesis directly and manipulate our expression; we had to multiply both sides by \(1 + x\) and then expand and 
simplify to show that the inequality still holds for \(n = k + 1\).

\QED

\section{Inequalities for Real Numbers}

For \(x,y \in \R\) the following holds: 

\[
    x < y + \varepsilon, \forall \varepsilon > 0 \iff x \le y
\]

and 

\[
    |x| < \varepsilon, \forall \varepsilon > 0 \iff x = 0.
\]

\textbf{Proof:}

For the first, if \(x > y\), then I can choose some epsilon only at some point where \(y + \varepsilon > x\); therefore, 
not all epsilons satisfy the inequality.

For the second one, if \(x \ne 0\), then \(|x| \ne 0 \land |x| > 0\); therefore, the only valid option is for 
\(x\) to be zero.

\QED

\subsection{Elementary Young Inequalities}

For all \(x,y \in \R\), the following holds:

\[
    xy \le \frac{1}{2}x^2 + \frac{1}{2}y^2.
\]

And this implies:

\[
    \sqrt{xy} \le \frac{1}{2}x + \frac{1}{2}y, \quad \forall x,y \ge 0
\]

and

\[
    xy \le \frac{\varepsilon}{2}x^2 + \frac{1}{2\varepsilon}y^2, \quad \forall x,y \in \R, \varepsilon > 0.
\]

\section{Inequalities With and Without Absolute Value}

\emph{Inequalities} are like normal \emph{equalities}, but with ranges of solutions 
instead of one particular solution. The inequalities without absolute value are straightforward 
to solve, just like equations, with the particularity that when dividing both sides by a negative number, 
the relation \(<, >, \ge, \le\) has to be converted to its opposite.

The case with absolute value is similar, but here we have to account for different cases depending on the 
expression inside the \(|\circ|\).

\textbf{Example:}

Solve \(|x + 1| \le 10 + |2x - 6|\).

\textbf{Main cases:}

\[
    F_1: x + 1 \ge 0 \implies x \ge -1
\]

\[
    F_2: x + 1 < 0 \implies x < -1
\]

\textbf{Subcases:}

\[
    F_{1,1}: 2x - 6 \ge 0  \implies x \ge 3
\]

\begin{align*}
    x + 1 &\ge 10 + 2x - 6 \\
    x &\le 5 \\
    \mathcal{L} &= [3; \infty )
\end{align*}

\[
    F_{1,2}: 2x - 6 < 0  \implies x < 3
\]

\begin{align*}
    x + 1 &\ge 10 - (2x - 6) \\
    x &\ge 5 \\
    \mathcal{L} &= (-\infty; 3]
\end{align*}

\[
    F_{2,1}: 2x - 6 \ge 0 \implies x \ge 3
\]

\begin{align*}
   -(x + 1 ) &\ge 10 + 2x - 6 \\
    x &\le \frac{-5}{3} \\
    \mathcal{L} &= \emptyset
\end{align*}

\[
    F_{2,2}: 2x - 6 < 0 \implies x < 3
\]

\begin{align*}
   -(x + 1 ) &\ge 10 - (2x - 6) \\
    x &\le 17 \\
    \mathcal{L} &= (- \infty, 3)
\end{align*}

\textbf{Union of the solutions:}

\[
    \mathcal{L} = 
    (-\infty; 3] \cup 
    [3; \infty ) \cup 
    \emptyset    \cup
    (- \infty, 3) = (- \infty, \infty)
\]


